% x86lean: the semantics paper (paper 1). TACAS 2027 regular track, LNCS format.
%
% Every number and every claim about another system carries a \src{...} comment naming the file in
% this repository it was taken from. The macro expands to nothing in the PDF, so the sources travel
% with the text and never print. Nothing is retyped from a message or a memory.
%
% Framing (ruled 2026-09-12, docs/TACAS-PRICING.md 2.1): a validated executable semantics. The Hoare
% logic is the second paper's claim; the measured tier law is a section here, not the headline.
%
% Status: DRAFT, every section written; no placeholder remains (the last, in Section 7, was replaced
% 2026-09-12, D222-D223). A revision pass has read every section against its markers, for length and
% register, and cited the paper's numbers to the manifest (docs/TACAS-PRICING.md 2.2, G7; D231-D235).
% This comment read "Every other section is a placeholder" until 2026-09-12, and "The revision pass is
% owed" until 2026-09-13.
\documentclass[runningheads]{llncs}
\usepackage[T1]{fontenc}
\usepackage{booktabs}
\usepackage{array}
\usepackage{amsmath}
\usepackage{url}
\usepackage[hidelinks]{hyperref}
\usepackage{microtype}
\newcommand{\src}[1]{}  % source-file marker; expands to nothing in the PDF
\newcommand{\lean}[1]{\texttt{#1}}
\newcommand{\placeholder}[1]{\par\noindent\fbox{\parbox{0.95\linewidth}{\small\emph{Not yet drafted.} #1}}\par}

\begin{document}

\title{x86lean: A Differentially Validated Executable Semantics\\of x86-64 User-Level Instructions in Lean~4}
\titlerunning{x86lean}
\author{Jason Hickey}
\authorrunning{J. Hickey}
\institute{Independent}

\maketitle

\begin{abstract}
x86lean is an executable semantics of x86-64 user-level instructions in Lean~4, built from public
sources and validated by differential testing against ACL2 x86isa. It covers fewer instructions than
the K semantics and has far weaker proof support than x86isa or K, and we say so first. What it contributes is a
treatment of validation in which the claims about a semantics, what it covers, how it was checked and
where the check has nothing to check against, are derived and checked rather than written, as far as the
gates reach, and the paper says where they stop. We describe a design for undefined bits that keeps one
definition for execution and proof, derives the undefined flags from the semantics instead of declaring
them, checks undefined registers against a declared rule, and states a case the model declines as a theorem; a
differential harness whose disagreement classes and wrong-model controls are explicit; the places where
the reference model is not the specification, including one where it contradicts itself; and a
measurement, on a sample, of what safety proofs over the semantics cost. We also report how often this
project's own ungated claims failed a re-check: in each of three samples, at least one time in three.
\src{Each claim is the subject of a section below and is cited there.}
\end{abstract}

% ---------------------------------------------------------------------------------------------
\section{Introduction}
\label{sec:intro}
% Register: docs/TACAS-PRICING.md 4a; docs/TACAS-G1-POSITIONING.md 1 and 3.

A formal semantics of a processor is used as a specification, so the first question about one is
whether it is right. The public semantics of x86-64 answer it by execution: x86isa is co-simulated
against a processor, and K was tested against thousands of instruction-level cases
(Section~\ref{sec:related}). This paper is about the second question, which a reader of any of them
also has to take on trust: are the claims made \emph{about} the semantics right? A coverage figure, a
list of what was validated against what, a statement of which bits are undefined, a count of
disagreements: each is ordinarily prose, and prose is not checked against the artifact it describes.

We begin with the concessions, because a reader comparing x86lean with its neighbours will reach them
first. It covers 158 mnemonics, where K's coverage target lists 774, and its proof support is a safety
theorem for each of six small routines, where x86isa proves full functional correctness of an
unbounded loop. Other formal x86-64 models exist beyond the three we compare against, and one of them,
libLISA, describes itself as the most extensive to date, with over 118{,}000 instruction groups
synthesized from a processor~\cite{liblisa}; Section~\ref{sec:related} says what we have and have not
read.
\src{docs/CLAIMS.tsv[coverage\_mnemonics, k\_mnemonics, proof\_sample\_routines, liblisa\_groups]; docs/TACAS-G1-POSITIONING.md, the scale and proof-support rows and 1d; docs/DECISIONS.md D220}
Its decoding is trusted to Intel XED, and its validation is against one executable model.

What x86lean contributes is a discipline for the claims around a semantics, applied throughout:
\begin{itemize}
\item an undefined-bit design in which one Lean definition serves execution and proof, the undefined flags
  are derived from the semantics rather than declared, undefined registers are checked from two sources,
  and a case the model declines is stated as a theorem (Section~\ref{sec:undef});
\item a differential harness whose classes of disagreement, units and wrong-model controls are explicit,
  a statement that the three models it could be compared against are two origins and not three, and the
  cases where the reference model is not the specification (Section~\ref{sec:validation});
\item coverage figures derived from two sources, and validation figures re-derived from their records,
  gated in CI against the prose that states them, this paper's manifest-cited figures among them, with a
  measurement of how often ungated claims in this project failed a re-check (Section~\ref{sec:gated});
\item a measurement, on a sample, of what safety proofs over the semantics cost, and a property the
  proof interface cannot state (Section~\ref{sec:proofs}).
\end{itemize}
Most of these ideas have public precedents, which Section~\ref{sec:related} cites. Their
combination is checkable in one clone of a public repository.

% ---------------------------------------------------------------------------------------------
\section{The Model}
\label{sec:model}
% Register: X86/State.lean Cpu; X86/Basic.lean GPR, Xmms; X86/Syntax.lean Instr; X86/Semantics.lean step,
% run; X86/Coverage.lean Tier, DecodeTrust; TRUSTBASE.md; docs/COVERAGE.md.

\paragraph{State and step.}
The machine state is a record: the sixteen general-purpose registers, the instruction pointer, the
arithmetic flags, memory, the FS and GS segment bases, sixteen 128-bit XMM registers, the undefined-bit
oracle, and a halt field that records why the model stopped, if it did.
\src{X86/State.lean structure Cpu; X86/Basic.lean GPR and Xmms}
The FS and GS bases are inputs: no instruction in the model writes them, so they are left out of the
differential record, and what is compared is the address they produce.
\src{X86/State.lean, the fsBase/gsBase docstring}
An instruction is an operation from the AST together with its encoded length, which the decoder
supplies. \lean{step} is a total function from an instruction and a state to a state. When an
instruction faults, or when the model declines a case, \lean{step} sets the halt field with a reason
rather than inventing a result; Section~\ref{sec:undef:refusal} gives an
example stated as a theorem. \lean{run} folds \lean{step} over a list of decoded instructions; the
model does not yet fetch and decode from memory.
\src{X86/Syntax.lean structure Instr; X86/Semantics.lean step and run}

\paragraph{Coverage as data.}
What the model claims is a table in Lean, not prose: one row per mnemonic, carrying its operand shapes,
a fidelity tier, how its decoding is trusted, and its undefined bits. A test module checks the table
against the AST in both directions, so a row cannot name a form the AST lacks and a form cannot lack a
row, and the published coverage document is regenerated from the table and compared in CI.
\src{X86/Coverage.lean module docstring; .github/workflows/ci.yml, the COVERAGE.md regeneration check}
The tier is \emph{exact} when every architected bit of the result and the flags is computed,
\emph{frame} when the set of changed state components is exact but some bits come from the oracle
because the manual leaves them undefined, and \emph{absent} when the model refuses. Of 158 mnemonics,
132 are exact, 26 are frame, and 0 are absent. The tier names are this project's coinage.
\src{X86/Coverage.lean inductive Tier; docs/COVERAGE.md:13; docs/CLAIMS.tsv[coverage\_mnemonics, tier\_exact, tier\_frame,
tier\_absent]}

\paragraph{Trust.}
Decoding is trusted, not proved. The AST is produced from Intel XED's structured output, and every one
of the 158 rows records its decoding as trusted to XED; a Lean decoder for the covered subset is later
work.
\src{TRUSTBASE.md, Decode trust; X86/Coverage.lean DecodeTrust; docs/CLAIMS.tsv[decode\_xed\_rows]}
Every theorem outside the test tree and a separately compiled native library depends on at most
\lean{propext}, \lean{Classical.choice} and \lean{Quot.sound}. The CI check is an allowlist of those
three names per declaration, not a denylist: on the Lean release we use, \lean{native\_decide}
introduces an axiom with a per-computation name rather than \lean{Lean.ofReduceBool}, so a denylist
keyed on the latter passes it silently, and the check is driven red with a planted
\lean{native\_decide} before it is trusted.
\src{TRUSTBASE.md, Axiom policy; scripts/axiom\_gate\_selftest.sh}
The test tree is outside that check, and it holds the six safety theorems of Section~\ref{sec:proofs} and
the theorems this paper cites from \lean{Tests}. No proof there uses \lean{native\_decide} or
\lean{bv\_decide}; a search establishes that at the published commit, and no gate enforces it.
\src{scripts/axiom\_gate.sh derives its modules from X86.lean's imports; Tests/Program.lean holds the six safety theorems;
git grep -n -E 'native\_decide|bv\_decide|ofReduceBool' -- 'Tests/*.lean' finds one line, a comment in Tests/Anchors.lean; docs/DECISIONS.md D248}
Two further boundaries are stated rather than modelled. The model assumes SSE is enabled, computing
vector results where hardware with \texttt{CR4.OSFXSR} clear would fault, and the reference model is
configured to match, a configuration that is itself checked. The LOCK prefix is recorded, and which
forms accept it is enforced, but atomicity is not verified: the semantics is single-threaded.
\src{TRUSTBASE.md, SSE availability and Atomicity; scripts/check\_driver\_cr4.py}

% ---------------------------------------------------------------------------------------------
\section{Undefined Behaviour}
\label{sec:undef}
% Register: docs/TACAS-G2-SECTION-DRAFT.md, docs/TACAS-G2-UNDEFINED-BITS.md,
% docs/TACAS-G1-POSITIONING.md 1c.6 and 1c.10, docs/DECISIONS.md D5 D6 D20 D52 D213 D226.

Throughout this paper the \emph{reference model} is ACL2 x86isa~\cite{goel-x86isa,x86isa-repo}, the
executable model the differential runs compare against, and the \emph{oracle} is the source of
undefined bits inside x86lean's own state. The two are unrelated.

The Intel manual~\cite{intel-sdm} leaves part of the processor state undefined after many
instructions, and it does so in more than one way. Three cases from the shift family show the range.
A \emph{flag} can be left undefined by a named clause: CF is undefined for ``SHL and SHR instructions
where the count is greater than or equal to the size (in bits) of the destination operand''.
\src{docs/DECISIONS.md D20; the wording re-read at the cited edition's PDF (325462-092US, SAL/SAR/SHL/SHR, Flags Affected) 2026-09-15,
where it carries ``(in bits)'', which this quotation dropped until then; X86/Flags.lean:194 has it; docs/DECISIONS.md D249}
A \emph{result} can be undefined: for SHLD, and in almost the same words for SHRD, ``if the count is
greater than the operand size, the result is undefined'', and the undefined part is the destination
together with all six arithmetic flags.
\src{docs/DECISIONS.md D52; both descriptions and both Flags Affected sections (``If the count is greater than the operand size, the
flags are undefined'') re-read at 325462-092US 2026-09-15, docs/DECISIONS.md D249}
And a case can look undefined without being so. An arithmetic right shift past the operand width has
no such clause and needs none: every vacated bit, and the last bit shifted out, is the sign bit. A
model that treats ``count at least the width'' as uniformly undefined invents undefinedness the manual
does not grant. x86lean gives SAR its own CF rule at those counts.
\src{X86/Flags.lean shiftFlags; docs/DECISIONS.md D20}
The model would have been self-consistent under either reading; it was the differential run of
Section~\ref{sec:validation} that confirmed this one, when a deliberately wrong copy of the model that
flips CF on a SAR case at such a count was contradicted by the reference model.
\src{docs/DECISIONS.md D20, the planted control on sar\_b9/6}

\subsection{Four Designs}
\label{sec:undef:designs}

Four public x86-64 semantics whose sources we have read, x86isa~\cite{goel-x86isa,x86isa-repo},
the K semantics~\cite{dasgupta-k-x86,k-x86-repo}, the Sail model translated from
x86isa~\cite{sail-x86-repo} and an Isabelle/HOL model in the Archive of Formal Proofs~\cite{afp-x86},
together with x86lean, answer the question ``what is an undefined bit?'' in five different ways
(Table~\ref{tab:undef}). Each row was read at the project's own source.
\src{docs/TACAS-G1-POSITIONING.md 1c.6 and 1c.10}

\begin{table}[t]
\caption{What an undefined bit is, in five public x86-64 semantics.}
\label{tab:undef}
\centering
\small
\begin{tabular}{@{}l>{\raggedright\arraybackslash}p{0.72\linewidth}@{}}
\toprule
system & an undefined bit is \\
\midrule
x86lean & a concrete bit read from an oracle carried in the machine state \\
ACL2 x86isa & in the logic, the value of \lean{create-undef}, a function introduced by
  \lean{encapsulate} with only a local witness, fed by a seed counter in the state; for execution, a
  concrete implementation attached under a trust tag \\
K x86-64 & a distinguished constant (\lean{undefMInt}, \lean{undefBool} and width-specific variants)
  written directly into the state \\
Sail x86 & Sail's built-in \lean{undefined}, reached through \lean{undef\_read\_logic}, together
  with a per-case mask of undefined flags passed to the flag writer \\
AFP Isabelle & a flag read with no defined semantics resolves to \lean{unknown\_flags i f}, where
  \lean{unknown\_flags :: string $\Rightarrow$ string $\Rightarrow$ bool} is fixed by a locale \\
\bottomrule
\end{tabular}
\end{table}
\src{x86isa: machine/register-readers-and-writers.lisp:1497; the trust tag in tools/execution/top.lisp.
K: semantics/x86-mint-wrapper.k:19. Sail: model/prelude.sail:74-82, model/shifts\_spec.sail,
model/rotate\_and\_shift.sail. All via docs/TACAS-G1-POSITIONING.md 1c.6 and 1c.10. AFP:
X86\_Semantics/X86\_InstructionSemantics.thy, locale unknowns and its text, read 2026-09-12 (G1 1d).}

These are different designs, not variations on one. x86isa's documentation states the property its
design is built for: because nothing is known about \lean{create-undef}, ``we wouldn't be able to prove
that a value obtained from \lean{undef} is equal (or not) to any other value \ldots\ an undefined value
is different from another undefined value, and also all the known values''~\cite{x86isa-repo}.
\src{docs/TACAS-G1-POSITIONING.md 1c.10, verbatim from create-undef's :long}
In K, 497 of the 3{,}064 per-instruction semantic files mention an undefined constant, and none of
them is a system instruction.
\src{docs/CLAIMS.tsv[k\_undef\_files, k\_perinst\_files]; docs/TACAS-G1-POSITIONING.md 1c.6: the denominator is the five directories K's README names as
holding per-instruction semantics; register 161, memory 227, immediate 109, system 0}
We make no claim about how K's or Sail's tools decide equality between two undefined values: we have
read both models' sources on that question and have run neither.
\src{docs/TACAS-G2-SECTION-DRAFT.md 5: measured at both sources, unrun against K}
The Isabelle/HOL model, of roughly 120 instructions, is by its own description ``purposefully
incomplete, but overapproximative'', and it deliberately avoids Isabelle's \lean{undefined}, ``since that
could be used to prove that the semantics of two undefined behaviors are equivalent''~\cite{afp-x86}.
Its unknown flag values are the results of a function fixed by a locale, applied to the instruction and
the flag name. Undefinedness as a parameter is therefore not new with us. What differs is the key: that
model's value is determined by the instruction and the flag, x86lean's by a position in a stream that
advances with each draw.
\src{docs/CLAIMS.tsv[afp\_instructions]; isa-afp.org/entries/X86\_Semantics.html and X86\_InstructionSemantics, locale unknowns, read
2026-09-12; docs/TACAS-G1-POSITIONING.md 1d; D220}
A sixth answer inverts the question. libLISA synthesizes instruction semantics by observing a processor,
and its model is CPU-specific: behaviour ``that is `undefined' is synthesized for the current
machine''~\cite{liblisa}. We have read its models only for the two rules of
Section~\ref{sec:validation:limits} on which x86lean and the reference model part, where they serve as a second source.
\src{docs/LIBLISA-HARDWARE-CHECK-2026-09-12.md; docs/DECISIONS.md D221}

x86lean's oracle is x86isa's seed discipline made concrete.
\src{PROVENANCE.md, the x86isa row; docs/TACAS-G6-RELATED-WORK.md 4}
What that buys is that the function which runs and the function which theorems are about are the same
Lean definition: \lean{step} reads undefined bits from the state, so it is an ordinary total function
that can be evaluated on a pre-state, and a theorem that places no hypothesis on the oracle holds for
every resolution of the undefined bits. x86isa keeps a logical story and an execution story apart; we
do not have to. What it costs is that undefinedness becomes a parameter the harness and the theorems
must both thread, and the three mechanisms below are what keep that parameter honest.

\subsection{The Draw Is Part of the Model}
\label{sec:undef:draw}

Which oracle bit lands in which field is fixed by the model. A shift whose masked count is non-zero
reads three bits, in the order CF, OF, AF; the logical group (AND, OR, XOR, TEST) reads one, for AF.
\src{X86/Semantics.lean, the .shift case; docs/DECISIONS.md D5; Tests/Nonvacuity.lean 3}
The rule that matters is that, once an instruction's branch is decided, its draws are unconditional: a
bit is drawn for a flag whether or not that flag is undefined at these operands. SAR at a count past the
width draws CF's bit and ignores it, because CF is the sign; SHL at a count of one draws OF's bit and
ignores it, because OF is defined there.
\src{X86/Flags.lean shiftFlags}
The number of bits drawn does still depend on operand values through the coarse branch: a masked count
of zero draws nothing, and BSF or BSR draw a further destination-width of bits only when the source is
zero.
\src{Tests/Nonvacuity.lean cursor\_shift\_zero\_spends\_nothing; X86/Semantics.lean bitScanStep;
docs/DECISIONS.md D213}
What the rule is meant to exclude is dependence on the oracle's own bits. If a branch inside a step
could read a bit drawn in that step, two runs under different oracles could take different branches,
and the derivation of Section~\ref{sec:undef:derived}, which compares two such runs, would report a
difference of control as a set of undefined fields.
\src{X86/Semantics.lean, the bsf/bsr and division comments}
We state the evidence at its actual strength. The rule is a design rule with example theorems that fix
the oracle cursor's position after particular instructions under two opposite oracles, and we have not
proved it for every instruction.
\src{Tests/Nonvacuity.lean cursor\_independent\_of\_bits; docs/DECISIONS.md D213}
It is checked case by case: the leak check of Section~\ref{sec:undef:derived} also requires the two runs
to leave the cursor in the same place, and over the 89{,}056 cases the differential compares the cursors
differ in 0. A model planted to draw one extra bit whenever its first drawn bit is set is caught on exactly
the cases that draw.
\src{docs/CURSOR-LEAK-CHECK-2026-09-13.md; docs/CLAIMS.tsv[cursor\_cases, cursor\_leaks]; X86/Serialize.lean
undefinedLeakedBy; Main.lean wrongDrawCountReadsDrawnBit; docs/DECISIONS.md D226}
The check cannot see a branch on a drawn bit that changes only the value a flag receives: a flag's
undefinedness is derived from those same two runs, with no second source to disagree.

\subsection{The Undefined Set Is Derived, Not Declared}
\label{sec:undef:derived}

The differential harness needs to know, for each case, which disagreements with the reference model are
explained by undefinedness. For flags it holds no list. \lean{undefinedFlags} runs the same
\lean{step} from the same pre-state under the all-zeros and the all-ones oracle and reports the flags
that differ.
\src{X86/Serialize.lean undefinedFlags; docs/DECISIONS.md D6}
A declared list would be a second statement of what the model leaves undefined, and it would go stale
the first time a form's rules changed. It would go stale in the permissive direction, because a list
that still names a flag the model now defines licenses the harness to excuse a real disagreement in it.
Derived, there is one notion of ``undefined here'', and it is the semantics'.

Registers need the opposite treatment, and the reason is instructive. BSF and BSR leave their
destination register undefined when the source is zero (``If the content of the source operand is 0,
the content of the destination operand is undefined''). Deriving the undefined registers from the same
two runs would disarm the check that matters most: an oracle bit reaching a register by mistake would
be re-read as ``that register is undefined here'' and filed as explained. So registers have two
sources, required to agree. \lean{declaredUndefGPRs} reads the instruction and the manual's rule;
\lean{undefinedRegs} observes which registers differ between the two runs. A register that moves
without being declared is a leak; a register that is declared and does not move is a false claim of
undefinedness, which would otherwise license the harness to explain away disagreements in that register
indefinitely.
\src{X86/Serialize.lean declaredUndefGPRs, undefinedRegs, and the P1 batch 14 note}

Everything outside those two channels must agree exactly between the two runs.
\lean{undefinedLeaked}, checked on every differential case, requires the vector registers, the
instruction pointer, the oracle cursor, the halt state and every watched memory window to be identical
under both oracles,
and the comparator will explain a disagreement as undefined only in a closed list of fields: the flags
and the general-purpose registers.
\src{X86/Serialize.lean undefinedLeaked; Main.lean undefinableFields}
The derivation has a precondition it does not itself check. It finds a field whose value is a copy of a
drawn bit, and it would miss a value that combined two drawn bits so as to agree under both constant
oracles. Every undefined value in the model is written as drawn; we established this by a syntactic
search of the model's sources, with a planted positive control, and not by proof.
\src{docs/DECISIONS.md D213 4}

\subsection{Declining a Case, as a Theorem}
\label{sec:undef:refusal}

SHLD and SHRD mask their count to five bits, so at a 16-bit operand the count can exceed the width and
the destination becomes undefined. This is not a corner: over the 82 pre-states of the harness at the
time, a count taken from CL exceeded 16 in 35.
\src{docs/DECISIONS.md D52, measured on the reference model before the constructor existed}
With a register destination, x86lean answers from the oracle, and the register channel above carries
it. With a memory destination there is no channel: \lean{undefinedLeaked} demands that the two runs
agree on every watched byte, and the comparator cannot explain a memory disagreement at all. Rather than
widen its strongest check as a side effect of one instruction family, the model declines exactly this
combination (a memory destination, a 16-bit operand, a masked count above 16) and halts with a stated
reason. Every other double shift is answered.
\src{X86/Syntax.lean dshiftMemUndefined; X86/Semantics.lean, the .dshift case}

The refusal is stated as a theorem,
\lean{step\_dshift\_mem\_\allowbreak undefined\_\allowbreak refuses}: for an immediate count and an effective address
without a LOCK prefix, the step leaves the state unchanged except for the halt reason. The guard in \lean{step} is on the masked count, which a count in CL
reaches as well.
\src{X86/Theorems.lean step\_dshift\_mem\_undefined\_refuses}
Each of the other four designs in Table~\ref{tab:undef} has a general means of answering such a case,
a constrained value, a constant, the built-in or a locale function; we have not read each model's rule for this
particular instruction, and so do not claim that any of them answers it. We claim no more for this than it is: one
instruction pair, presented as a demonstrated discipline rather than as coverage. It is, however,
self-enforcing. The reference model computes a result for this case, so a test vector there would be
scored as a one-sided refusal, not as a pass; the gap cannot be filled by accident, and because it is a
theorem it cannot be quietly forgotten either.
\src{docs/DECISIONS.md D52}

% ---------------------------------------------------------------------------------------------
\section{Differential Validation}
\label{sec:validation}

\subsection{The Harness}
\label{sec:validation:harness}
% Register: Main.lean classify, compareRecs, bothRefused, driveWrong; X86/Serialize.lean Cpu.render;
% docs/DIFFERENTIAL-P2-BATCH22.md; docs/CLAIMS.tsv diff\_*.

A differential case is one instruction, run from one pre-state. x86lean and the reference model each
emit a record of the post-state in one format: the general-purpose and vector registers, the
instruction pointer, the flags, whether the instruction was refused, and a set of watched memory
windows.
\src{X86/Serialize.lean Cpu.render; scripts/run\_differential.sh}
The records are compared field by field. A case with no disagreeing field is \emph{matched}; a case in
which both models refuse is matched on the refusal alone, since neither claims a post-state. Each
disagreeing field is then classified. It is \emph{explained} if the field is one the comparator may
excuse (a flag or a general-purpose register) and x86lean's derived undefined set for that case
contains it; it is a \emph{declared divergence} if an entry names that vector and that field; it is a
\emph{refusal} if only one side refused, a \emph{harness} failure if a record lacks the field, and a
\emph{specification} disagreement otherwise. Only the first two classes are excused.
\src{Main.lean classify, bothRefused, undefinableFields, knownDivergences}

The summary counts two different things, and they must be read apart: cases, and disagreeing fields.
The current record covers 1{,}012 vectors at 88 pre-states, 89{,}056 cases, of which 68{,}524 are
matched. The remaining cases carry 29{,}435 field disagreements explained by the undefined set and 171
declared divergences. 0 disagreements are unexplained, and in 0 cases did an undefined bit leak.
\src{docs/DIFFERENTIAL-P2-BATCH22.md:4-5; docs/CLAIMS.tsv[diff\_vectors, diff\_prestates, diff\_cases,
diff\_matched, diff\_explained, diff\_diverged, diff\_unexplained, diff\_leaks]. Units: Main.lean
compareRecs counts cases, matched and leaks per case, and the other three per field}
A re-run with the reference model pinned at ACL2 commit c8897a34 reproduced every counter: 89{,}056 cases,
68{,}524 matched, 29{,}435 explained, 171 declared divergences, 0 unexplained and 0 leaks. Both leak counts
were taken before the leak check compared the oracle cursor; the cursor's own count is the separate record of
Section~\ref{sec:undef:draw}.
\src{docs/REFERENCE-PIN-RUN-2026-09-12.md; scripts/oracle\_revision.txt; docs/CLAIMS.tsv[pin\_reference\_model\_short, pin\_cases,
pin\_matched, pin\_explained, pin\_diverged, pin\_unexplained, pin\_leaks]}
A comparator that never reports a disagreement is indistinguishable from one that cannot. The harness
therefore runs deliberately wrong copies of x86lean against the correct one, with an empty divergence
list, and requires each to produce a specification or refusal disagreement in the field its defect
touches.
\src{Main.lean driveWrong: compareRecs [] good bad; hits exclude undefined-region and harness}

\subsection{Two Origins, Not Three Witnesses}
\label{sec:validation:origins}
% Register: docs/TACAS-G1-POSITIONING.md 1c.5, 1c.7; PROVENANCE.md; Main.lean knownDivergences.

The three public x86-64 models we position x86lean against could easily be counted as three witnesses. They are
two origins. The Sail x86 model is, in its repository's own description, automatically translated from the
ACL2 model, and its validation guide co-simulates it against K's single-instruction tests; agreement
with it would test the translation, not supply an independent semantics.
\src{docs/TACAS-G1-POSITIONING.md 1c.5 and 1c.7}
Our use of the two origins is also unequal, and we state it. Every differential run in this paper is
against x86isa alone. K has never been executed by this project: its rule files are read, for coverage
and to arbitrate a disagreement between x86lean and x86isa, as in the packed-shift case of
Section~\ref{sec:validation:limits}, where each declared divergence cites the K rule that agrees with
x86lean. Sail has not been used in any differential record.
\src{PROVENANCE.md, the K and Sail rows; docs/TACAS-G1-POSITIONING.md, the executable row and 1c.7;
Main.lean knownDivergences, the source field of every entry}
The validation claims below are therefore claims of agreement with one executable model, arbitrated
where the two disagree by the manual and by reading the other model's rules, and for the two rules behind
every declared divergence by a model synthesized on processors (Section~\ref{sec:validation:limits}). That arbiter is fallible
too: comparing K's semantics with semantics synthesized on five processors, libLISA judged K incorrect for
between 18 and 28 instruction variants per machine~\cite{liblisa}.
\src{docs/CLAIMS.tsv[liblisa\_k\_incorrect\_min, liblisa\_k\_incorrect\_max]; liblisa.nl/publications/liblisa-oopsla24/05-results, section 5.2.2 comparison table, row
``Dasgupta et al.\ incorrect'': 28 24 18 18 18, read 2026-09-12; docs/TACAS-G1-POSITIONING.md 1d}
Those verdicts cannot have reached the rules we cite. On none of its machines do libLISA's semantics contain
an instruction led by a 66 prefix, which leaves out the non-VEX forms of SSE instructions, as that work says;
every K rule behind our declared divergences is such a form, and a verdict on a K rule needs a synthesized
semantics to set against it.
\src{liblisa.nl/publications/liblisa-oopsla24/05-results, section 5.2.2, ``Out-of-scope'', read 2026-09-12;
docs/LIBLISA-HARDWARE-CHECK-2026-09-12.md, the prefix census over all five machines; docs/DECISIONS.md D221}

\subsection{Where the Reference Model Is Not the Specification}
\label{sec:validation:limits}
% Register: docs/TACAS-G2-SECTION-DRAFT.md 4; docs/DECISIONS.md D91 D108 D115; Main.lean
% knownDivergences; docs/DIFFERENTIAL-P2-BATCH13.md; docs/P2-ROSTER.md.

A paper whose argument is that every claim is checked must also report where the check has nothing to
check against. Four findings fall there.

\paragraph{A rule with no second source.}
MOVDQA faults on a memory operand that is not 16-byte aligned. Asked by execution, the reference model
does not: it executes both an unaligned load and an unaligned store, while an aligned MOVDQA, an
unaligned MOVDQU and a control instruction behave as expected.
\src{docs/DECISIONS.md D91 1}
x86lean faults, as the manual says, and no unaligned MOVDQA vector is shipped, since every one would be
scored as a one-sided refusal against a model that is right. The rule is therefore carried by theorem
rather than by run: an unaligned load faults, an unaligned store faults, and an unaligned MOVDQU at the
same address does not, which is the theorem that shows the alignment flag is doing the discriminating.
\src{docs/DECISIONS.md D91 2; X86/Theorems.lean vload\_unaligned\_faults, vstore\_unaligned\_faults,
vload\_unaligned\_movdqu\_runs}
This is the weakest claim in the model: it rests on our reading of the manual and on a proof that the
model implements that reading, and nothing else.
Agreement with a reference model that does not implement a rule is not evidence about the rule.

\paragraph{A reference model that keeps bits the manual clears.}
Moving a general-purpose register into an XMM register (MOVD, MOVQ) clears the destination's upper bits in
the manual and in K's rules; the reference model keeps them and writes only the low 32 or 64. x86lean
follows the manual, and the two vectors are declared divergences with K's rule as their independent
source. They are most of the declared count: 163 of the 171. A wrong model planted to merge rather than
clear was already one of the batch's controls, and it fired against x86lean; the reference model
implements the same mistake. libLISA's semantics agree with the manual on the VEX-encoded forms, on each
of its five machines.
\src{docs/DECISIONS.md D93, sections 1 and 2 (the planted arm, and K's movd\_xmm\_r32.k and movq\_xmm\_r64.k);
Main.lean knownDivergences, movd\_to\_x and movq\_to\_x; docs/DECISIONS.md D108, the first run's record
``unexplained=8 oracle-divergence=163'', taken while those two were the only declared entries, and 171 once the
eight shifts were declared; docs/LIBLISA-HARDWARE-CHECK-2026-09-12.md, the vmovd and vmovq verdicts; D221; docs/DECISIONS.md D248}

\paragraph{A reference model that disagrees with itself.}
On the packed shifts (PSLLW, PSRAD and their siblings) with a count in an XMM register, the first
differential run returned 8 unexplained disagreements in 78{,}584 cases, all at one pre-state: the only
one whose count register had a small low quadword and a non-zero high one. x86lean shifted by the low
quadword; the reference model saturated.
\src{docs/CLAIMS.tsv[d108\_unexplained, d108\_cases]; docs/DECISIONS.md D108 1}
The manual reads the count from the low 64 bits, and so does K: its saturation test extracts exactly
those bits of the source. More decisively, the reference model's own memory-count form of the same
instruction, given the same 128-bit count, returns the manual's answer.
\src{docs/DECISIONS.md D108 3, citing K's psllw\_xmm\_xmm.k}
A processor agrees too, on a sibling encoding. libLISA's semantics, synthesized on five machines, contain
none of these instructions with the 66 prefix our vectors use, but they contain the VEX-encoded forms, for
which the manual gives the same count rule. On each machine, the result of each of the eight shifts depends
on the low 64 bits of the count register and on none of the upper 64: 40 verdicts, none failing, beside a
control instruction whose upper-half read is visible and plants that each turn a verdict red. The two
encodings differ above bit 127, which this reading does not examine.
\src{docs/LIBLISA-HARDWARE-CHECK-2026-09-12.md; scripts/liblisa\_hardware\_check.py;
docs/CLAIMS.tsv[liblisa\_shift\_verdicts\_pass]; docs/DECISIONS.md D221}
A model that disagrees with itself across two operand shapes of one instruction has a defect, not a
reading. We record each of the eight as a known divergence naming one vector and one field, with its
independent source; the cases keep running, and an entry that stops diverging fails the run, so a fix
upstream is noticed rather than silently absorbed.
\src{Main.lean knownDivergences; docs/DECISIONS.md D108 4}
The published record of that batch reads the same 78{,}584 cases with 0 unexplained and 0 leaks, the
eight now counted as divergences and the matched count unchanged.
\src{docs/DIFFERENTIAL-P2-BATCH11.md:138-139, the eleventh differential record, which is the seat's batch 13 and D108's; docs/CLAIMS.tsv[p2\_shift\_record\_cases, p2\_shift\_record\_unexplained];
matched 58,052 and explained 29,435 in both records, divergences 163 -> 171. This sentence read 81,752 until 2026-09-15, the count of
DIFFERENTIAL-P2-BATCH13.md, which is a later batch's record with 36 more vectors (the two counters that file's header names); docs/DECISIONS.md D249}
A divergence entry has a cost, and we state it. It excuses one field of one vector at every
pre-state, so a genuine defect of x86lean in that field would now be absorbed by it. Two things stand
against that, and neither is the entry itself. The saturation rule is checked by a theorem on a witness
value at every lane width, chosen so that a model returning zero, or confusing a logical with an
arithmetic shift, fails it. And the deliberately wrong models of these instructions are compared
against x86lean itself with an empty divergence list, where the reference model's defect is
irrelevant, so no entry can weaken them.
\src{docs/DECISIONS.md D108 4; Tests/Coverage.lean vshift\_saturates\_rather\_than\_wrapping;
Main.lean driveWrong, compareRecs [] good bad}

\paragraph{The ceiling of the method.}
Differential validation needs a reference model that runs the instruction, and for much of what
x86lean does not yet cover, ours does not. We measure demand with a census of the instructions in
public Debian binaries; in its assembly class, six codec columns counted apart from compiler output,
401{,}653 instructions are forms the model does not cover.
\src{docs/DEMAND-CENSUS.md, the corpus recipe; docs/P2-ROSTER.md header; docs/CLAIMS.tsv[p2\_gap\_total]}
Of those, 162 are scalar instructions in the census row for the LOCK vocabulary, still counted as
uncovered although the vocabulary itself has landed, and the rest are vector instructions.
Asking the reference model to execute one form of each vector mnemonic, at the operand class where the
demand actually occurs, it executes the forms carrying 139{,}854 of them and refuses the forms carrying
169{,}877; on 241, in three mnemonics, it neither advances nor refuses; and 91{,}519 have not yet been
asked at their own class. The five parts sum to the total.
\src{docs/P2-ROSTER.md, the additions table and the BY (mnemonic, BUCKET) table; docs/CLAIMS.tsv[p2\_addition\_lock,
p2\_bucket\_executes, p2\_bucket\_refuses, p2\_bucket\_stalls, p2\_bucket\_unasked, p2\_ceiling\_partition\_residue];
162 + 139854 + 169877 + 241 + 91519 = 401653}
Among the forms it has been asked, the reference model therefore refuses more than it executes; the
unasked share is larger than that difference, so the comparison does not yet extend to the whole gap.
The refused instructions cannot be differentially validated against it at any price, however the work
is ordered. K has
semantic rules for many of the refused mnemonics, but it executes assembly source rather than machine
bytes and this project has never run it, so whether it can serve as the second source there is open.
\src{docs/DECISIONS.md D115 3 for the finding (its table is dated; the live figures are the roster's);
docs/P2-ROSTER.md, the K operand shapes column; docs/TACAS-G1-POSITIONING.md 1c.8 and its executable row}
For those instructions a model built from the manual has, for now, no second source, and a paper about
validation has to say so.

% ---------------------------------------------------------------------------------------------
\section{Gated Claims}
\label{sec:gated}
% Register: scripts/claimed_forms.py; docs/COVERAGE.md:5; docs/CLAIMS.tsv and scripts/check_claims.py;
% docs/DECISIONS.md D201 D202 D213 D214; docs/TACAS-G6-RELATED-WORK.md 2.

A semantics makes two kinds of claim: what its rules say, and what it covers and how it was checked.
Theorems police the first. The second is usually prose, and prose is a second register that nothing
compares with the first. We treat such claims as outputs.

\paragraph{Coverage.}
The coverage figures are computed from two sources that are not derived from each other. One is each
differential vector's own assembly text, parsed into rows of the scalar roster derived from K; the
other is each roster row, synthesised into a canonical instance and assembled by clang. A vector may
claim a row only if the two assemble to the same opcode, so a mis-parse cannot survive; rows whose
instances assemble to identical bytes are reported as spellings of one machine form, which is why the
figures count both rows and distinct forms.
\src{scripts/claimed\_forms.py module docstring}
The published figures are 158 mnemonics in 1{,}012 differentially tested vectors, covering 500 of the 525
roster rows, which are 351 of the 374 distinct machine forms those rows describe; CI fails if the prose
that states them disagrees with the derivation.
\src{docs/COVERAGE.md:5 and README.md:56-57, checked against the derivation by scripts/claimed\_forms.py --check in CI;
this paper's copy is checked against docs/CLAIMS.tsv[coverage\_mnemonics, coverage\_forms, coverage\_rows\_covered,
coverage\_rows\_total, coverage\_forms\_covered, coverage\_forms\_distinct], which read COVERAGE.md at the pin (D231)}
Numbers outside the coverage table are entered in a manifest with the command that re-derives each
and the commit it is published at; CI runs every command and fails on a difference. The manifest is
not yet complete, and it records which claims it holds rather than implying the rest. The commit is pinned
because a measurement is a claim about a commit: re-derived at the head, a count of decisions would go
red every time a decision was recorded.
\src{docs/CLAIMS.tsv header; scripts/check\_claims.py}

\paragraph{Why.}
We did not adopt this from principle; we measured the alternative on ourselves. Re-driving a table of
prior-art claims this project had written nine days earlier, each marked as verified at its source,
three of nine failed, and the files they described had not changed, so the claims were wrong when they
were written.
\src{docs/DECISIONS.md D201 1; docs/CLAIMS.tsv[recheck\_d201\_failed, recheck\_d201\_total]}
A bibliography table headed as fetched from the registrar disagreed with the record its own DOI returns
in three of eight rows.
\src{docs/DECISIONS.md D214; docs/CLAIMS.tsv[recheck\_d214\_failed, recheck\_d214\_total]}
Seven headline counts in a planning document were all correct and none stated its population;
re-derived with the obvious commands, three of the seven came out different.
\src{docs/DECISIONS.md D202; docs/CLAIMS.tsv[recheck\_d202\_failed, recheck\_d202\_total]}
And the technical draft behind Section~\ref{sec:undef}, written from decision records, contained two
false sentences that surfaced only when the prose was checked against the step function.
\src{docs/DECISIONS.md D213}
Writing this paper's related work added two more cases: attributions to prior work that the project
had recorded from the literature and never re-read did not survive reading the paper cited, one crediting
a frame rule to a paper whose improvement removes it and one describing as hardware validation what the
paper reports as validation against test suites and simulators.
\src{docs/DECISIONS.md D222 and D223. Not counted as a sample, and not in the rate below: the set they were
drawn from was not fixed before the reading, and D223's own record counts three corrections that evening
where this paragraph once counted two failures of three; docs/DECISIONS.md D235}
In the three samples with a denominator, 3 of 9, 3 of 8 and 3 of 7, claims that no check read did
not survive a re-check at least one time in three: wrong in the first two samples, and in the third not
reproducible from what they stated. In the first sample, where it could be tested, the failures were
present when the claims were written rather than acquired later.
Re-checking more often does not answer that; only a mechanical check does.
\src{docs/CLAIMS.tsv[recheck\_d201\_failed, recheck\_d201\_total, recheck\_d214\_failed, recheck\_d214\_total,
recheck\_d202\_failed, recheck\_d202\_total]; the draft of Section~3 and the related-work cases have no denominator
and are not in the rate (D235)}

\paragraph{What this is not.}
Deriving a claim and failing the build on disagreement is ordinary engineering, and it is not novel in
kind. K derives its own list of unsupported instructions with a script rather than declaring it.
\src{docs/TACAS-G6-RELATED-WORK.md 2: scripts/find\_unsupported\_insttructions.pl}
What we add is the gate on the prose, and its application to a semantics' coverage and validation
claims, where a referee can clone the repository and watch it run.
The paper is held to the same rule. Each number it quotes from the manifest carries, in its source,
a citation naming the rows it quotes, and the build fails if a cited value is absent from the sentence
it annotates, if a citation names no row, or if a row that says it appears in the paper is cited
nowhere. A number with no citation is not checked by this, and we do not claim that it is.
\src{scripts/check\_claims.py check\_prose and its selftest arms; docs/DECISIONS.md D216}

% ---------------------------------------------------------------------------------------------
\section{Proofs Over the Semantics}
\label{sec:proofs}
% Register: docs/P2-PROOF-INTERFACE.md; X86/Program.lean; Tests/Program.lean; docs/TACAS-PRICING.md 1;
% docs/TACAS-G1-POSITIONING.md, the proof-support row.

This paper's claim is the validated semantics; this section reports what it costs to prove safety
properties over it, measured on a small sample, and it is a sample and not a benchmark.

\paragraph{A driver that follows control.}
The differential harness drives \lean{run}, a fold of \lean{step} over a list of instructions, and that
is correct for single instructions. It cannot express a loop. On a three-instruction countdown the
branch is taken in the semantics, since the instruction pointer after the fold is the loop label, but
the fold proceeds to the next list element regardless, and returns a well-formed state that has not
halted, the answer after one pass.
\src{docs/P2-PROOF-INTERFACE.md, the finding before either design question: ecx 2, rip at the label,
no halt}
\lean{runP} executes a program, a table of decoded instructions at addresses, for a given amount of
fuel, fetching the instruction at the current instruction pointer at each step and halting when the
pointer leaves the table. A general invariant rule over \lean{runP} yields the safety theorems below.
The countdown's is accompanied by theorems that pin its concrete trace, driven red first, because a
safety property holds vacuously of a program that halts at once.
\src{X86/Program.lean runP, runP\_invariant, runP\_safe; Tests/Program.lean, the nonvacuity theorems}

\paragraph{Two tiers of property, and their cost.}
A property that constrains which state components may change, such as ``this routine writes no
memory'', is proved over every start state and every amount of fuel without reasoning about labels. On
the countdown loop and on a four-instruction scan, each such theorem is 14 lines, which on inspection
is a fixed base of about eleven lines plus about one line per instruction.
\src{docs/P2-PROOF-INTERFACE.md, problem 3 section 1; the equal counts are a coincidence of line
wrapping; counted by scripts/proof\_lines.py; docs/CLAIMS.tsv[proof\_lines\_countdown, proof\_lines\_scan]}
A property that depends on position, such as ``every store lands inside the destination buffer'', needs
an invariant indexed by label. Across four routines with 2, 3, 5 and 7 labels, the proofs are 34, 50, 57
and 76 lines. A least-squares line through the four is about 22 lines plus 7.6 per label, or about 175
lines at twenty labels; the routine furthest from it, by 5 lines, is the one whose store depends on a
bounds check, and four points do not show whether the per-label cost rises with the label count.
\src{scripts/label\_fit.py (a label is an entry of the invariant's table, cross-checked against the program's
instructions; lines by scripts/proof\_lines.py); docs/CLAIMS.tsv[fit\_labels\_prologue, fit\_labels\_guarded,
fit\_labels\_fill, fit\_labels\_memcpy, proof\_lines\_prologue, proof\_lines\_guarded, proof\_lines\_fill,
proof\_lines\_memcpy, fit\_intercept, fit\_slope, fit\_at20, fit\_max\_residual]; docs/DECISIONS.md D234: the
guarded routine was recorded as 4 labels, and the earlier fit of 19 plus 7.5 to 8.1 per label rested on it}
Three successive rounds of interface work on the routine with five labels brought its per-label cost,
measured above that routine's own fixed base, from about 25 lines to 12, 9.6 and 7.6; the third was pre-registered with a refutation threshold, which
it did not reach, and it added no lemma at all: its saving came from deleting proof steps that
definitional unfolding already discharged and from one macro.
\src{docs/CLAIMS.tsv[fit\_labels\_fill, proof\_round\_1, proof\_round\_2, proof\_round\_3, proof\_round\_4]; docs/P2-PROOF-INTERFACE.md, the runP\_code point and round 4 (problem 1, the fill routine; this sentence said "the first of these routines", which in the list above is prologue, D234)}
We read this as the tier deciding the cost more than the library does: a factor of about eight between
the two tiers, against the roughly threefold reduction the three rounds achieved within one. We do
not claim a floor. The residue per label, unpacking the invariant, naming the instruction's effect
theorem and re-establishing the invariant, is proof content rather than repetition, which is why we
expect the remedy to be a tactic or a verification-condition generator rather than more lemmas.
\src{docs/P2-PROOF-INTERFACE.md round 4: the residue argument, stated as an argument}

\paragraph{What cannot be stated.}
One property in the sample cannot be expressed as an invariant of a single run: that a scan reads only
inside its buffer. Memory reads are total and leave no trace in the state, so a wild read is
unobservable to any single-run invariant. The honest form is non-interference over two runs whose
memories agree on the buffer, and the interface's combinators are all single-run.
\src{docs/P2-PROOF-INTERFACE.md, problem 3 section 2}

\paragraph{The sample.}
Six routines, four of them with two to seven labels, written by one author against one interface, are a sample.
Against x86isa's eight verified program families, including full functional correctness for an
unbounded copy loop, this is a far weaker proof development (Section~\ref{sec:related}). A benchmark of such properties, and a program logic proved sound over this
semantics, are the subject of later work.
\src{docs/TACAS-PRICING.md 2.1 (the ruling) and G4; docs/TACAS-G1-POSITIONING.md 1c.2; docs/CLAIMS.tsv[proof\_sample\_routines]}

% ---------------------------------------------------------------------------------------------
\section{Related Work}
\label{sec:related}
% Register: docs/TACAS-G1-POSITIONING.md (the table and 1b, 1c), docs/TACAS-G6-RELATED-WORK.md (1, 4, 5,
% 6), docs/DECISIONS.md D201 D214. Every bibliography record fetched by DOI (D214).

\paragraph{x86-64 semantics.}
We compare against three of the x86-64 semantics whose sources we have read: two independent of each other,
and a third derived from one of them.
ACL2 x86isa~\cite{goel-x86isa,x86isa-repo} specifies ``400+ opcodes executing in Intel's 64-bit mode of
operation'' and runs ``co-simulations against an actual x86 processor for model validation''.
\src{docs/CLAIMS.tsv[x86isa\_opcodes]; docs/TACAS-G1-POSITIONING.md 1b, verbatim from arXiv 1705.01225, both phrases in its section 1
(re-read at the PDF 2026-09-15, D249), published as EPTCS 249}
Its proof development is far stronger than ours: eight verified program families carry 559 theorem
forms over a separate proof library, and one of them proves full functional correctness and freedom
from faults for a copy loop of unbounded length.
\src{docs/CLAIMS.tsv[x86isa\_theorem\_forms]; docs/TACAS-G1-POSITIONING.md 1c.2: dataCopy/dataCopy.lisp copyData-is-correct; 559 excludes the
603 forms of the utilities library}
The K semantics~\cite{dasgupta-k-x86,k-x86-repo} covers ``3155 instruction variants, corresponding to
774 mnemonics'', and ``is fully executable and has been tested against more than 7,000
instruction-level test cases and the GCC torture test suite''; its prover establishes functional
post-conditions for ten programs, including a loop.
\src{docs/CLAIMS.tsv[k\_variants, k\_mnemonics, k\_tests]; docs/TACAS-G1-POSITIONING.md 1b verbatim from K's README; 1c.4 for the kprove specifications}
K executes assembly source rather than machine code: its loader parses assembler directives, and its
opcodes are mnemonics.
\src{docs/TACAS-G1-POSITIONING.md 1c.8}
The Sail x86 model~\cite{sail-x86-repo} is translated automatically from x86isa, carries none of its
theorems, and is validated by co-simulation against K's single-instruction tests, so agreement with it
tests the translation rather than supplying a third semantics.
\src{docs/TACAS-G1-POSITIONING.md 1c.5 and 1c.7}
Against these, x86lean covers 158 mnemonics where K's target lists 774; x86isa's figure is stated in
opcodes, a different unit, and Sail's model publishes no count. Its proof support, a safety theorem for
each of six small routines, is far weaker than x86isa's or K's; the Sail model carries no theorems.
\src{docs/CLAIMS.tsv[coverage\_mnemonics, k\_mnemonics, proof\_sample\_routines]; docs/TACAS-G1-POSITIONING.md, the scale and proof-support rows and 1d;
158 is derived by scripts/claimed\_forms.py into COVERAGE.md, and this copy is checked against the row (D231); docs/DECISIONS.md D220}
Other formal x86-64 models exist outside this comparison, and the largest changes what validation
means. libLISA discovers instructions and synthesizes their semantics by executing them on a CPU,
producing ``the most extensive formal x86-64 model to date, with over 118 000 different instruction
groups'', compares models across five machines, and reports exposing bugs in handwritten
models~\cite{liblisa}. Where x86lean asks whether a model built from the manual agrees with another
model, libLISA takes the processor as ground truth. The two are complementary rather than competing.
We have used libLISA's models for the packed-shift rule on which x86lean and the reference model part
(Section~\ref{sec:validation:limits}); a comparison of x86lean against them beyond those rules is the most
direct second source this work still lacks, and for the SSE forms it would first need their scope extended.
\src{docs/CLAIMS.tsv[liblisa\_groups]; docs/LIBLISA-HARDWARE-CHECK-2026-09-12.md; docs/DECISIONS.md D221; the quoted ``118 000'' keeps libLISA's own spacing, which the prose check reads since D235}
A public Isabelle/HOL entry in the Archive of
Formal Proofs gives over-approximative semantics for roughly 120 x86-64 assembly instructions, with
symbolic execution of basic blocks and a parser for disassembler output~\cite{afp-x86}. Roessle et al.\ extract a machine model
with small-step semantics for 1{,}625 instruction variants from the semantics stratified synthesis learned,
and generate equivalence proofs between decompiled machine code and big-step semantics~\cite{roessle-bigstep};
Verbeek et al.\ lift stripped binaries to a representation that provably over-approximates their execution
paths, exportable to Isabelle/HOL, and take sound instruction semantics as an assumption~\cite{verbeek-lifting}.
We have read these papers, not their artifacts, and do not position x86lean against them: the first counts
instruction variants rather than mnemonics, and the second is about lifting, not an instruction set.
\src{docs/CLAIMS.tsv[afp\_instructions, roessle\_instructions]; abstracts via api.semanticscholar.org for DOIs 10.1145/3293880.3294102 and 10.1145/3519939.3523702,
and 10.1145/3689723, and isa-afp.org/entries/X86\_Semantics.html, read 2026-09-12; docs/DECISIONS.md D220;
the two papers read at the authors' PDFs 2026-09-14 (ssrg.ece.vt.edu/papers/cpp2019.pdf and pldi22.pdf, titles and authors matched to
their DOI records): CPP 2019 section 1, ``a machine learned and formally tested x86-64 machine model containing 1625 instructions variants (IVs)'',
and section 2, ``Strata demonstrates trustworthy semantics for 692 instructions, which through generalization arguments expands to 1625 IVs'';
PLDI 2022 section 1, ``a representation that contains an overapproximation of all possible execution paths of the binary'' (abstract) and
``We also assume the existence of sound instruction semantics''; docs/DECISIONS.md D246}
What x86lean adds is the subject of Sections~\ref{sec:undef}--\ref{sec:gated}: an undefined-bit design
whose limits are stated, one of them as a theorem, a differential harness whose disagreement classes and wrong-model
controls are explicit, and coverage and validation claims that are derived and checked rather than
written.

\paragraph{Techniques.}
Most of x86lean's techniques have public precedents, and we cite them where they originate. A total step function over
one machine state is the shape of all three models above. The closest precedent in Lean~4 is
LNSym~\cite{lnsym-repo}, a model of an Arm instruction set with co-simulation conformance testing and a
benchmarks target. It models a different instruction set, so we cite it for the method and do not
compare coverage with it.
\src{docs/TACAS-G6-RELATED-WORK.md 1 row 7 (VERIFIED at its README) and the LNSym ruling in 6}
Validation by differential testing against executable models and hardware is established practice:
K was tested against more than 7{,}000 instruction-level cases~\cite{dasgupta-k-x86}, x86isa is
co-simulated against a processor~\cite{goel-x86isa}, and stratified synthesis learned the x86-64
instruction set automatically~\cite{heule-stratified}.
\src{docs/CLAIMS.tsv[k\_tests]; docs/TACAS-G6-RELATED-WORK.md 4, the differential and co-simulation rows;
each attribution re-read at its source 2026-09-13 (D236): the PLDI 2019 abstract via api.semanticscholar.org for DOI
10.1145/3314221.3314601, ``has been tested against more than 7,000 instruction-level test cases''; arXiv:1705.01225's
section 1, ``run co-simulations against an actual x86 processor for model validation'' (this marker said ``abstract'' until
2026-09-15: the phrase is not in the abstract, re-read at DataCite and the PDF, D249); the PLDI 2016 title, ``Automatically
Learning the x86-64 Instruction Set''}
The check against a processor can itself
be a proof: the model extracted from those learned semantics was tested in Isabelle/HOL by lemmas, each
proving the step function's result for a pre-state whose post-state a processor had computed~\cite{roessle-bigstep}.
\src{CPP 2019 section 6, read at the authors' PDF 2026-09-14: ``For each IV, we create Isabelle/HOL test lemmas'';
``The post-execution states are computed using dynamically generated binaries executed on live x86-64 hardware (Skylake
architecture)''; 886 variants tested, two failing; docs/DECISIONS.md D246}
Using XED as a trusted decoder is not a clean contrast with x86isa either: 186 of the 3{,}192 entries in
x86isa's opcode maps were generated from XED's data files.
\src{docs/CLAIMS.tsv[x86isa\_xed\_entries, x86isa\_inst\_entries]; docs/TACAS-G1-POSITIONING.md 1c.3 and 1c.9}
Specifying which parts of a machine state code may change, with the rest held as a frame, has a direct
precedent for machine code: Myreen and Gordon's Hoare logic over realistically modelled machine code
writes specifications in the style of separation logic and derives its frame rule~\cite{myreen-hoare-machine-code},
and it was instantiated to models of x86, ARM and PowerPC~\cite{myreen-multiple-architectures}. Its later
form derives one theorem per instruction by evaluating the instruction-set model once and composes those
theorems without evaluating the model again~\cite{myreen-decompilation-improved}. Our program proofs
likewise cite an instruction's effect theorem rather than unfolding the step function, though ours are
proved by hand and theirs are derived automatically.
\src{author PDFs and the author's bibtex (cse.chalmers.se/\textasciitilde myreen), read 2026-09-12: TACAS 2007
section 3.2, ``only the part initially satisfying P has been changed'' and ``The following frame-rule, similar
to that of separation logic, easily follows''; FMCAD 2008 section III, the frame rule derived from the triple; FMCAD 2012 sections II and III, Phase 1 and Fig.~1;
Tests/Program.lean cites step\_* effect theorems and never step itself; docs/TACAS-G6-RELATED-WORK.md 5;
docs/DECISIONS.md D222}
The Sail models of ARMv8-A, RISC-V and CHERI-MIPS were validated against executable references and test
suites rather than against a processor: the ARM model ran ARM's Architecture Validation Suite,
where 24 of 15{,}400 tests passed on ARM's own specification and failed on the translation, and the RISC-V
model's execution traces were compared against the Spike reference simulator~\cite{armstrong-isa}.
\src{docs/CLAIMS.tsv[armstrong\_avs\_asl\_only, armstrong\_avs\_tests], which read this reading's record, docs/DECISIONS.md D223 (D235); Armstrong et al., POPL 2019, section 7 (Validation), read at cl.cam.ac.uk/\textasciitilde pes20/sail/sail-popl2019.pdf
2026-09-12: ``Of those 15 400 tests, currently 24 (0.15\%) pass on the ASL model but not on the Sail model'';
``compare the trace outputs of the Sail model and a version of Spike''; docs/DECISIONS.md D223}
That a model should state its own fidelity is not our idea either: the Isabelle/HOL model above describes
its semantics as over-approximative~\cite{afp-x86}. What the fidelity tier of
Section~\ref{sec:model} adds is that the statement is made per form, as data that is checked against the AST.
\src{isa-afp.org/entries/X86\_Semantics.html, ``purposefully incomplete, but overapproximative''; X86/Coverage.lean inductive Tier and
the two-direction table test cited in Section~\ref{sec:model}; docs/DECISIONS.md D223. This sentence also cited the lifting work
(verbeek-lifting) until 2026-09-14; its over-approximation is of a binary's execution paths, and it assumes the instruction
semantics (PLDI 2022 section 1), so it states no model's fidelity: docs/DECISIONS.md D246}

% ---------------------------------------------------------------------------------------------
\section{Limitations and Conclusion}
\label{sec:conclusion}

The limitations are the ones stated where they arise. x86lean is small, and its proofs are few. Its
decoding is trusted. Its validation is agreement with one executable model, arbitrated by the manual and
by reading a second, and for two rules by a model synthesized on processors; K has not been run, and no
hardware co-simulation has been. Of the remaining instructions it has
been asked, the reference model refuses more than it executes, and for those a model built from the
manual currently has no second source. A rule the reference model does not implement, such as MOVDQA's
alignment fault, is carried by theorem alone. The proof-cost measurements are a sample of six routines. And the
gate on quoted figures checks the figures that cite the manifest, not every number in the text.

The contribution is not a larger or more strongly proved semantics. It is a semantics whose claims about
itself are outputs: derived from the artifact, checked against a second source where one exists, and
stated as absent where one does not. The measurement that motivates it is modest and, we think,
general: in this project, claims about its artifacts and its sources that nothing checked did not survive
a re-check at least one time in three in each of the three samples we could count.

% ---------------------------------------------------------------------------------------------
% Mandatory for ETAPS 2027 papers: placed just before the references, outside the page limit
% (docs/TACAS-PRICING.md 3a, read at etaps.org/2027/cfp/).
\section*{Data Availability}
% Written on the case-study arm (docs/TACAS-PRICING.md 3a). A research-paper submission is double-blind and
% needs this anonymized. The submission commit is not yet fixed.
The model, its proofs, the differential harness and every record cited here are public at
\url{https://github.com/jyh/x86lean} under the Apache-2.0 licence. Every figure this paper quotes from
the manifest \texttt{docs/CLAIMS.tsv} names the commit it is pinned at and the command that re-derives
it, and \texttt{python3 scripts/\allowbreak check\_\allowbreak claims.py} re-derives all of them and checks the paper's text
against them. The Lean toolchain is pinned by \texttt{lean-toolchain}; the reference model is ACL2
x86isa at ACL2 commit \texttt{c8897a34}, pinned by \texttt{scripts/oracle\_revision.txt} and fetched
by \texttt{scripts/setup\_oracle.sh}.
\src{LICENSE; lean-toolchain; scripts/oracle\_revision.txt; docs/CLAIMS.tsv[pin\_reference\_model\_short, proof\_sample\_routines]}

\bibliographystyle{splncs04}
\bibliography{x86lean-semantics}

\end{document}
